\chapter{What a layer draws}\label{ch:elements}
\begin{goals}
\textbf{In this chapter} you choose what each layer draws (poles, great circles, contours, the mean and
its cone, eigenvectors) and set how they draw in the layer's properties.
\end{goals}

A layer is more than its points. Open it in the tree and it unfolds into everything it can draw, each with
a checkbox: OpenStereo's nested checkboxes, kept because they work. What is ticked is drawn; how it is drawn
is set once in the layer's properties, and then you go back to looking at the net.

\section{The nested checkboxes}\label{sec:checkboxes}

\shot{elements-tree}{\code{tocher} opened into its elements: contours (Kamb, 3\,$\sigma$), the mean with its
95\,\% cone, and the eigenvector group partly on (V1 with its great circle, and V3). The statistics tab
below gives the numbers.}

Click the small arrow beside a layer to open it. Each row is an \term{element}; its checkbox shows or hides
it, and is an undo step (\ui{show / hide elements}). Some rows carry their setting at the right:
\ui{Kamb 3σ} beside contours, \ui{95 \%} beside the mean.

\term[eigenvectors!group]{The eigenvectors} are a group with a checkbox of its own. Ticking it switches all six
of its children on, unticking switches them off, and when only some are on its box shows a dash: it is
\emph{partly} on. Open the group to tick its children one by one.

What a layer can draw depends on its dataset's kind, and a new layer starts with the simplest thing on:

\begin{center}\small
\begin{tabularx}{\linewidth}{@{}>{\raggedright}p{0.2\linewidth}L>{\raggedright\arraybackslash}p{0.22\linewidth}@{}}\toprule
\sffamily\bfseries kind & \sffamily\bfseries elements & \sffamily\bfseries on at first \\\midrule
planes & poles, great circles, contours, mean + cone, eigenvectors (V1–V3, each a point and a great circle), labels & poles \\
lines & points, contours, mean + cone, eigenvectors, labels & points \\
faults & great circles, slip lines, P and T axes, right dihedra, labels & great circles, slip lines \\
small circles & circles, axes, labels & circles \\
\bottomrule\end{tabularx}
\end{center}

\begin{why}
Two hundred great circles bury a net. Poles show the same planes without hiding each other, so a new
layer shows poles (OpenStereo does the same). Faults are the exception: their great circles and slip
lines are the point of them (\cref{ch:faults}).
\end{why}

The console reaches the same switches: \code{show contours}, \code{hide poles}, \code{show v1}
(\cref{sec:commands}).

\section{Poles, points and great circles}

\term{Poles} are drawn as dots, in the layer's colour and symbol size; a selected row's pole is violet and
larger, on every diagram that draws it. \term{Points} are the same for lines. \term{Great circles} draw
each plane in full, in a thin line of the layer's colour, violet when selected. A small-circle dataset
draws its \term{circles}, and \term{axes} adds a point at each circle's axis.

Clicking a pole or a point picks its row; \cref{ch:select} covers picking and the other selection tools.

\section{Contours}\label{sec:contours}

Contours show the density of the poles (or lines) on the sphere. Two methods:

\begin{description}
\item[Kamb] \parencite{kamb1959}, the default. Measurements are counted inside a circle on the sphere whose
size is set so that the count expected from a uniform distribution is $k$ standard deviations of it:
$E = k\sigma$, with $k = 3$ unless you change it. The contours are in $\sigma$ above that expectation, at
2, 4, 6, 8, 10 and 12\,$\sigma$.
\item[Fisher kernel] each measurement contributes a smooth Fisher bump, and the result is expressed in
\term[multiples of uniform density]{multiples of uniform density} (MUD): contours at 2, 4, 6 and 8 times
what a uniform spread would give.
\end{description}

Contours are drawn under the layer's circles and points, in its colour. They need at least three readable
measurements.

\begin{caution}
A net draws \emph{one} set of contours: those of the first layer, in draw order, that has contours on.
Tick contours on a second layer of the same net and the net says, in its corner,
\ui{contours: one set per net — “west” isn’t contoured; the contours are “east”’s}. Contour the second
dataset on a net of its own (\menu{Diagram>New net}, then remove the first layer there).
\end{caution}

The contour levels themselves are a setting with no dialog yet: write a \code{"levels"} list in the layer's
\file{.json} under \file{/project} (\cref{sec:project-write}), for example
\code{"contours": \{"on": true, "method": "kamb", "sigma": 3, "levels": [2, 4]\}}. Set it back to
\code{null} for the defaults.

\section{Mean and cone}\label{sec:mean}

\term{Mean + cone} draws the layer's Fisher mean \parencite{fisher1953} as an open ring, and around it the
confidence cone at 95\,\% (or 90 or 99\,\%: properties). Poles and lines are treated as \emph{axes}: a pole
and its opposite are the same direction, so the poles of steep planes dipping both ways, near the rim on
opposite sides of the net, are one cluster with one mean, not two halves that cancel. The numbers ($\kappa$, $\alpha_{95}$, the mean as a plane and as a
line) are on the statistics tab (\cref{ch:stats}). It needs two readable measurements.

\begin{note}
A Fisher mean describes a cluster. On a girdle, such as the Tocher sample's, it lands somewhere
meaningless and its cone is wide; the eigenvectors below describe a girdle.
\end{note}

\section{Eigenvectors}\label{sec:eigen}

The orientation tensor of the poles (or lines) has three eigenvectors, \term{V1}, \term{V2} and \term{V3},
from the largest eigenvalue to the smallest. Each can be drawn two ways:

\begin{itemize}
\item \ui{V1}, \ui{V2}, \ui{V3}: the eigenvector as a dark point;
\item \ui{V1 great circle} and the others: the plane perpendicular to it, as a dashed great circle.
\end{itemize}

For a cluster, V1 is its centre. For a girdle, V3 is the axis of the girdle and \ui{V3 great circle} is
the best-fit girdle itself: on a set of folded bedding poles, that is the $\pi$ circle and V3 the fold axis.
Eigenvectors need three readable measurements; the eigenvalues and the fabric shape parameters are on the
statistics tab (\cref{ch:stats}).

\section{Labels}\label{sec:labels}

Every kind has a \ui{labels} element: a small label beside each measurement's point (the pole of a
plane, the line itself, a fault's slickenline), in the layer's colour, violet when the row is selected.
The properties dialog's \ui{labels from} chooses what it says: the row number (the default), or any
column, such as a station name. Labels are drawn whether or not the points are, so a net of great
circles can be labelled too, and they go into figures with the rest of the net.

\begin{tip}
On 200 measurements labels are a blur. Select the few you want to name, make them a dataset
(\cref{sec:select-dataset}), and label that layer.
\end{tip}

\section{Properties}\label{sec:properties}

\shot[0.8\linewidth]{layer-props}{A layer's properties: its colour and symbol size, and the settings of its contours,
mean cone and labels.}

Double-click a layer in the tree, or choose \ui{Properties…} from its menu. The dialog's first line says
what it is reading: \ui{tocher · planes · 200 of 200 rows read}. Unreadable rows are counted here
(\cref{sec:data-read}).

\begin{center}\small
\begin{tabularx}{\linewidth}{@{}>{\raggedright}p{0.24\linewidth}L>{\raggedright\arraybackslash}p{0.16\linewidth}@{}}\toprule
\sffamily\bfseries setting & \sffamily\bfseries meaning & \sffamily\bfseries default \\\midrule
\ui{colour} & the layer's colour: its points, circles, contours and cone & the next palette colour \\
\ui{symbol size} & the size of its points, 1 to 12 & 4 \\
\ui{contours} & \ui{Kamb (σ)} or \ui{Fisher kernel (MUD)} & Kamb \\
\ui{Kamb k (E = kσ)} & the $k$ that sets Kamb's counting circle, 1 to 6 & 3 \\
\ui{cone confidence} & the mean's cone: 90, 95 or 99\,\% & 95\,\% \\
\ui{labels from} & the row number, or a column (see \cref{sec:labels}) & row number \\
\bottomrule\end{tabularx}
\end{center}

Only the settings the layer's kind has are shown: a fault layer has no contours or cone. \ui{Apply} sets
them; \ui{Cancel} or \key{Esc} leaves the layer as it was. The settings are the layer's, not the data's: a
view of the same data on another net keeps its own.

\begin{tryit}
\begin{enumerate}
\item Open the Tocher sample and open \code{tocher} in the tree. Tick \ui{contours}: Kamb at 3\,$\sigma$.
\item Double-click the layer, set \ui{contours} to \ui{Fisher kernel (MUD)}, and \ui{Apply}. The tree
now says \ui{Fisher} beside contours, and the contours change shape.
\item Open \ui{eigenvectors} and tick \ui{V3} and \ui{V3 great circle}. The dashed circle runs along the
Tocher girdle, and V3 is its axis.
\item Tick the \ui{eigenvectors} group itself: all six come on, and the dash in its box becomes a tick.
\end{enumerate}
\end{tryit}
